\scanpage{113}
% archive foliation: 114
\margin{\fbox{Conférence à Nice}}
\uline{Groupes de Barsotti--Tate}
\fbox{$p$ fixé}

1) fppf sur $S$

$G$ gpe de B.T si
\[
\begin{array}{ll}
a)&pG=G\\
b)&G\text{ de $p$-torsion}\\
c)&{}_pG=\operatorname{Ker}p.\operatorname{id}_G\quad\text{fini loc. libre}
\end{array}
\]
Alors les ${}_{p^n}G=G(n)$ le sont aussi. Si $G(1)$ de rang $p^d$, alors $G(n)$ de rang $p^{nd}$ ; on a $G=\underrightarrow{\lim}_nG(n)$

Exemple important : $A$ schéma abélien,
\[
G(n)={}_{p^n}A\qquad G={}_{p^\infty}A
\]
\deleted{\unsure{2)}} Th. de Serre--Tate
\[
S_o\subset S\qquad A\longmapsto(A_o,{}_{p^\infty}A,\varphi)
\]
\margin{\deleted{\unsure{S...}} de car. $p$}
une équivalence de catégories

2) Groupe formel $\bar G$ associé à un groupe de B.T.

3) \uline{Cristaux}\qquad Puissances divisées sur un idéal $J$
\[
x^{(n)}\quad(n\geq0)\quad\left(\frac{x^n}{n!}\right)
\]
\margin{$x^{(0)}=1,\ x^{(1)}=x\qquad x^{(n)}\in J$}
\[
\left\{\begin{array}{l}
(x+y)^{(n)}=\displaystyle\sum_{i+j=n}x^{(i)}x^{(j)}\\[1em]
(\lambda x)^{(n)}=\lambda^n x^{(n)}\\[2em]
\ldots
\end{array}\right.
\]
\margin{$J^{(n)}$ \inserted{idéal} engendré par les monômes de poids $\geq n$ en les $x^{(i)}$ ($x\in J$)}

\deleted{\uuline{Site cristallin absolu}\quad$\{$\quad$\unsure{x}$}

puiss. div. nilpotentes : $J^{(n)}=0$ si $n$ grand

\rule{2cm}{.3pt} quasi nilpotentes : $J^{(n)}$ de torsion pour $n$ grand

i.e. $(n-1)!J^{(n)}=0$

exponentielle
\[
\exp(x)=\sum_{n\geq0}x^{(n)}\qquad\text{OK. si p.d. nilp.}
\]
\[
\log(1+x)=\sum_{n\geq1}(-1)^{n-1}(n-1)!x^{(n)}\qquad\text{OK.}\quad\rule{1cm}{.3pt}\quad\text{quasi-nilp.}
\]
\margin{\unsure{début de} Théorie cohomologique \unsure{développée par} Berthelot (et Illusie)}

\uline{Site cristallin} \deleted{\unsure{naturel|relatif}} \inserted{(absolu)}\quad\deleted{\unsure{cris}}

\unsure{nilpot. au sens Berthelot} (\unsure{si} $p$ nilp sur $S$)

On exige que les p.d. de $p$ soient les habituelles $\ldots$
